\documentclass[11pt,a4paper]{article}
\usepackage{fontspec,amsmath,amssymb,amsthm,microtype}
\setmainfont{lmroman10-regular.otf}[BoldFont=lmroman10-bold.otf,ItalicFont=lmroman10-italic.otf,BoldItalicFont=lmroman10-bolditalic.otf]
\usepackage{xcolor}
\usepackage[most]{tcolorbox}
\usepackage{mathrsfs}
\providecommand{\lt}{<}
\providecommand{\gt}{>}
\setmonofont{DejaVuSansMono.ttf}[BoldFont=DejaVuSansMono-Bold.ttf,ItalicFont=DejaVuSansMono-Oblique.ttf,BoldItalicFont=DejaVuSansMono-BoldOblique.ttf,Scale=0.86]
\newfontfamily\leanfallback{STIXTwoMath-Regular.otf}[Scale=0.86]
\definecolor{leanrule}{gray}{0.62}
\definecolor{leanback}{gray}{0.962}
\definecolor{leantitle}{gray}{0.22}
\definecolor{leanmuted}{gray}{0.38}
\newtcolorbox{leanbox}[1]{enhanced,breakable,colback=leanback,colframe=leanrule,
  boxrule=0.4pt,arc=0.6mm,left=2.4mm,right=2.4mm,top=1.8mm,bottom=1.8mm,
  before skip=0.8em,after skip=1em,fonttitle=\small\bfseries,coltitle=leantitle,
  colbacktitle=leanback,titlerule=0.3pt,title={#1}}
\newlength{\leancharwidth}
\newenvironment{leancode}{\par\smallskip\begingroup\ttfamily\footnotesize\raggedright
  \settowidth{\leancharwidth}{M}\setlength{\parskip}{0pt}\setlength{\parindent}{0pt}
  \frenchspacing}{\par\endgroup\smallskip}
\newcommand{\leanline}[2]{\hangindent=\dimexpr#1\leancharwidth+4\leancharwidth\relax
  \hangafter=1\noindent\hspace*{#1\leancharwidth}#2\par}
\newcommand{\leansource}[1]{\par{\footnotesize\color{leanmuted}#1}\par}
\newcommand{\leanlabel}[1]{\par\smallskip\noindent{\small\bfseries #1}\par}
\usepackage[margin=25mm]{geometry}
\usepackage[unicode,hidelinks]{hyperref}
\hypersetup{pdfauthor={LangevinLean Project}}
\setlength{\emergencystretch}{3em}
\setlength{\parskip}{0.4em}
\setlength{\parindent}{0pt}
\allowdisplaybreaks[2]
\raggedbottom
\title{\texorpdfstring{Preserving actual dynamics under an independent source extension}{Preserving actual dynamics under an independent source extension}}
\author{LangevinLean Project}
\date{}
\begin{document}\maketitle
Draft mathematical exposition; these notes may contain mistakes.


Appending an independent coordinate to the probability space changes neither the Langevin dynamics nor their laws. Measure-preserving pullback commutes with the true drift, Bochner time integral, Euler recursion and stationary observable evaluation.



Earlier results used here: \cite[Section~3.14]{LangevinLean2026Kinetic}; \cite[Section~3.16]{LangevinLean2026ULASC}.

\section{\texorpdfstring{Lifting genuine solutions and observables}{Lifting genuine solutions and observables}}


\textbf{Claim.} Pulling a Langevin solution and its additive noise through the first projection of an independent product gives the same integral equation on the extended space. Euler endpoints and stationary evaluations lift in the same way, and their laws are unchanged.



\textbf{Proof.} The first projection is measure-preserving because the added measure is a probability. Its \(L^2\) pullback is a linear isometry. The actual Lipschitz drift commutes with this pullback by comparing representatives. Bounded linear pullback commutes with the Bochner integral; induction gives the Euler identity. Exact map composition and measure preservation give the law and observable-evaluation identities.

\begin{leanbox}{Lean: \texttt{langevin\_\allowbreak{}integral\_\allowbreak{}sourceLift}, \texttt{additiveEulerChain\_\allowbreak{}sourceLift}, \texttt{sourceLift\_\allowbreak{}law} and 1 more}
\leansource{\texttt{langevin\_\allowbreak{}integral\_\allowbreak{}sourceLift} \quad WassersteinMixing/\allowbreak{}SemigroupLift.lean, lines 28–44}
\begin{leancode}
\leanline{0}{theorem\ langevin\_integral\_sourceLift\ (g\ :\ E\ →\ E)\ (L\ :\ ℝ≥0)\ (hg\ :\ LipschitzWith\ L\ g)}
\leanline{4}{(initial\ :\ Lp\ E\ 2\ P)\ (B\ Y\ :\ ℝ\ →\ Lp\ E\ 2\ P)\ (hY\ :\ Continuous\ Y)}
\leanline{4}{(T\ :\ ℝ)\ (heq\ :\ ∀\ t\ ∈\ Icc\ 0\ T,\ Y\ t\ -\ initial\ =}
\leanline{6}{-(∫\ s\ in\ 0..t,\ ULA.lpDrift\ P\ g\ hg\ (Y\ s))\ +\ B\ t)\ :}
\leanline{4}{∀\ t\ ∈\ Icc\ 0\ T,}
\leanline{6}{KineticLangevin.sourcePullback\ φ\ hφ\ (Y\ t)\ -}
\leanline{10}{KineticLangevin.sourcePullback\ φ\ hφ\ initial\ =}
\leanline{8}{-(∫\ s\ in\ 0..t,\ ULA.lpDrift\ Q\ g\ hg\ (KineticLangevin.sourcePullback\ φ\ hφ\ (Y\ s)))\ +}
\leanline{10}{KineticLangevin.sourcePullback\ φ\ hφ\ (B\ t)}
\end{leancode}
\leansource{\texttt{additiveEulerChain\_\allowbreak{}sourceLift} \quad WassersteinMixing/\allowbreak{}SemigroupLift.lean, lines 47–58}
\begin{leancode}
\leanline{0}{theorem\ additiveEulerChain\_sourceLift\ (g\ :\ E\ →\ E)\ (L\ :\ ℝ≥0)\ (hg\ :\ LipschitzWith\ L\ g)}
\leanline{4}{(step\ :\ ℝ)\ (initial\ :\ Lp\ E\ 2\ P)\ (B\ :\ ℝ\ →\ Lp\ E\ 2\ P)\ (n\ :\ ℕ)\ :}
\leanline{4}{KineticLangevin.sourcePullback\ φ\ hφ}
\leanline{6}{(ULA.additiveEulerChain\ (ULA.lpDrift\ P\ g\ hg)\ step\ initial\ B\ n)\ =}
\leanline{4}{ULA.additiveEulerChain\ (ULA.lpDrift\ Q\ g\ hg)\ step}
\leanline{6}{(KineticLangevin.sourcePullback\ φ\ hφ\ initial)}
\leanline{6}{(fun\ t\ ↦\ KineticLangevin.sourcePullback\ φ\ hφ\ (B\ t))\ n}
\end{leancode}
\leansource{\texttt{sourceLift\_\allowbreak{}law} \quad WassersteinMixing/\allowbreak{}SemigroupLift.lean, lines 65–67}
\begin{leancode}
\leanline{0}{theorem\ sourceLift\_law\ (X\ :\ Lp\ E\ 2\ P)\ :}
\leanline{4}{Q.map\ (KineticLangevin.sourcePullback\ φ\ hφ\ X)\ =\ P.map\ X}
\end{leancode}
\leansource{\texttt{stationaryEvaluation\_\allowbreak{}sourceLift} \quad WassersteinMixing/\allowbreak{}SemigroupLift.lean, lines 72–89}
\begin{leancode}
\leanline{0}{theorem\ stationaryEvaluation\_sourceLift\ (μ\ :\ Measure\ E)\ (X\ :\ Lp\ E\ 2\ P)}
\leanline{4}{(hX\ :\ P.map\ X\ =\ μ)\ (f\ :\ Lp\ F\ 2\ μ)\ :}
\leanline{4}{StationaryObservables.evaluation\ ({\leanfallback 𝕜}\ :=\ {\leanfallback 𝕜})\ μ\ Q}
\leanline{6}{(KineticLangevin.sourcePullback\ φ\ hφ\ X)\ ((sourceLift\_law\ φ\ hφ\ X).trans\ hX)\ f\ =}
\leanline{6}{Lp.compMeasurePreservingₗᵢ\ {\leanfallback 𝕜}\ φ\ hφ}
\leanline{8}{(StationaryObservables.evaluation\ ({\leanfallback 𝕜}\ :=\ {\leanfallback 𝕜})\ μ\ P\ X\ hX\ f)}
\end{leancode}
\leanlabel{How the Lean proof follows the argument}
\texttt{langevin\_\allowbreak{}integral\_\allowbreak{}source\allowbreak{}Lift} transports the actual integral equation through the pullback isometry. \texttt{additive\allowbreak{}Euler\allowbreak{}Chain\_\allowbreak{}source\allowbreak{}Lift} inducts on its recursion. \texttt{source\allowbreak{}Lift\_\allowbreak{}law} and \texttt{stationary\allowbreak{}Evaluation\_\allowbreak{}source\allowbreak{}Lift} use exact measure-preserving composition and representative equality.
\end{leanbox}


Earlier results used here: \cite[Section~3.13.2]{LangevinLean2026ULASC}; \cite[Section~3.13.5]{LangevinLean2026ULASC}; \cite[Section~3.16.1]{LangevinLean2026ULASC}; \cite[Section~3.16.3]{LangevinLean2026ULASC}; \cite[Section~3.14.1]{LangevinLean2026Kinetic}; \cite[Section~3.14.2]{LangevinLean2026Kinetic}.

\begin{thebibliography}{ULA-LSI}
\bibitem[ULA-SC]{LangevinLean2026ULASC}
LangevinLean Project. \emph{Unadjusted Langevin under Strong Convexity}. LL-01, 2026. Draft report, content revision d34ead06987c. \url{https://odakyildiz.github.io/LangevinLean/docs/pdf/UnadjustedLangevin/FullProof.pdf}.
\bibitem[Kinetic]{LangevinLean2026Kinetic}
LangevinLean Project. \emph{Kinetic Langevin Monte Carlo}. LL-06, 2026. Draft report, content revision b2fee42a9593. \url{https://odakyildiz.github.io/LangevinLean/docs/pdf/KineticLangevin/FullProof.pdf}.
\end{thebibliography}

\end{document}
